\section{相关工作}

\subsection{机器人模仿学习与少样本操作}

行为克隆通过监督学习拟合示范中的观测到动作映射。Robomimic~\cite{mandlekar2021robomimic} 系统研究了离线人类示范上的策略结构与训练选择。ACT~\cite{zhao2023act} 通过动作分块和 Transformer 学习精细双臂操作，Diffusion Policy~\cite{chi2023diffusion} 则将多模态动作序列建模为条件去噪过程。上述方法提升了策略表达能力，但不能直接解决少量示范未覆盖协同节奏的问题。本文不提出新的基础策略架构，而是研究一种可与 ACT、BC-RNN 或 Diffusion Policy 组合的数据生成与训练机制。

\subsection{自动示范生成}

MimicGen~\cite{mandlekar2023mimicgen} 将源示范划分为物体中心的片段，通过空间变换、拼接和仿真执行生成新轨迹。DexMimicGen~\cite{jiang2024dexmimicgen} 将该思想扩展到双臂和灵巧操作，引入逐臂子任务分割、协同子任务同步以及顺序约束。本文与其互补：DexMimicGen主要解决“如何在新场景中合成可执行双臂轨迹”，而本文关注“在动作几何基本不变时，左右臂相对时序如何决定成败，以及如何把这种信息用于少样本学习”。

\subsection{人类视频与跨具身数据}

EgoMimic~\cite{kareer2024egomimic} 将第一视角人手数据与机器人数据联合训练，说明人类具身数据可以扩展机器人学习的数据来源。不过，人手与机器人在运动学、视觉外观和动作分布上存在差异；更重要的是，视觉或运动学上合理的轨迹不一定满足机器人执行时的接触时序。本文当前版本从已重定向的机器人轨迹出发研究物理协同边界，未来可以将该模块接在人类视频重定向之后，作为迁移轨迹的物理筛选和时序修正环节。

\subsection{时序增强与双臂协同}

常见时序增强包括统一加速、减速、帧丢弃和随机动作延迟。统一时间缩放保持两臂相对相位，难以检验协同依赖；直接冻结或删除 delta 动作又会改变累计末端位姿，造成混杂。本文使用后缀延迟：在等待之后完整重放原动作后缀，使单臂内部动作次序和动作内容保持不变，从而更有针对性地研究相对时序。
